\section{$\minilang$ Grammar}
\label{sec:grammar}
This section presents the \emph{syntax} of $\minilang$,
a language of numbers and strings.
The following notation may seem unusual
because most programmers write programs in the
\emph{concrete syntax} of a language.
The concrete syntax of language specifies how
humans write programs in the language.
Type systems are often written over the
\emph{abstract syntax} of the language to
reflect that expressions in the language are
\emph{abstract syntax trees} (ASTs)
or more simply called \emph{terms}.
ASTs are mathematical-like expressions
that represent a composition of \emph{nodes}.
Each AST has the form: \\*
$operator(operand_1, operand_2, \ldots, operand_n)$,
where each operand is an AST, the operator
is a root node of these ASTs, and $n \geq 0$.
AST nodes are operators that take a specified
number of operands.
An operator can take in zero operands; in this case,
the parentheses after the operator are typically not written.
For instance, in the AST \code{add(6, 1)},
the node \code{6} could have been written as
\code{6()},
and \code{add(6, 1)} is equivalent to \code{add(6(), 1())}.
Only the abstract syntax is needed to reason about
a language.

The abstract and concrete syntax of a language is defined with a
\emph{formal grammar} that specifies \emph{production rules} on how AST
nodes are constructed.
Each production rule has the form
``\mbox{$A \bnfdef B_1 \bnfalt B_2 \bnfalt \ldots \bnfalt B_n$}''.
A \emph{non-terminal symbol} $A$ denotes a set of ASTs specified
by a production rule, where $A$ is on the left-hand side
of the $\bnfdef$.
Each $B_i$ denotes a category of AST nodes
or a \emph{terminal symbol} denoting a specific AST node.
The production rule specifies that the $A$ symbol denotes
the category of nodes that is the union of nodes
in $B_1, B_2, \ldots, B_n$.
Figure~\ref{fig:grammar} shows the grammar of $\minilang$.

\begin{figure}[H]
\[
\begin{array}{llrl@{\hspace{1.5cm}}l}
\mathit{Category}   & \mathit{Item}   & & \mathit{Abstract} & \mathit{Concrete} \\ \\
\mathit{Expression} & e & \bnfdef & x                              & x \\
                    &   & \bnfalt & \code{num[$n$]}                & n \\
                    &   & \bnfalt & \code{str[$s$]}                & \code{'$s$'} \\
                    &   & \bnfalt & \code{+($e_1$; $e_2$)}         & \code{$e_1$ + $e_2$} \\
                    &   & \bnfalt & \code{\^{}($e_1$; $e_2$)}      & \code{$e_1$ \^{} $e_2$} \\
                    &   & \bnfalt & \code{let($x$; $e_1$; $e_2$)}  & \code{let $x$ be $e_1$ in $e_2$}
\end{array}
\]
\caption{Grammar of $\minilang$}
\label{fig:grammar}
\end{figure}


The only non-terminal symbol in $\minilang$ is $e$,
which denotes the set of ASTs that are expressions.
Expressions can be the following:
\begin{enumerate}
\item Numbers: \code{num[$n$]},
where $n$ is a sequence of digits.
\item Strings: \code{str[$s$]},
where $s$ is a sequence of characters.
\item Variable names: $x$
\item The sum of two subexpressions: \code{+($e_1$, $e_2$)}.
\item The string concatenation of two subexpressions:
\code{\^{}($e_1$, $e_2$)}.
\item A let expression where the subexpression
$e_2$ is evaluated in a context that has variable
$x$ bound to the value of $e_1$. 
\end{enumerate}

Example expressions are shown in Figure~\ref{fig:ex_expressions}.
For the remaining sections in this paper, only the abstract
syntax will be shown.

\begin{figure}[H]
\begin{center}
\begin{tabular}{l|l}
\textit{Abstract} & \textit{Concrete} \\ \hline
\code{+(num[5]; +(num[4]; num[3]))}   & \code{5 + 4 + 3} \\
\\

\code{\^{}(str[john]; \^{}(x; str[doe]))} & \code{'john' \^{} x \^{} 'doe'} \\
\\

\code{let(hours; num[24]; +(hours; num[33])}    & \code{let hours be 24 in hours+33} \\
\end{tabular}
\end{center}
\caption{Example $\minilang$ Expressions}
\label{fig:ex_expressions} 
\end{figure}

